\begin{explanationbox}
	\textbf{\textcolor{CExplanation}{一句话直觉}}

	任意角的正弦和余弦的平方和固定为1，就像单位圆上点的横纵坐标平方和恒等于半径平方。

	\medskip
	\textbf{\textcolor{CExplanation}{核心拆解}}
	\begin{description}[style=nextline, leftmargin=2.8em, labelsep=0.5em, itemsep=0.45em]
		\item[\textbf{要点一（平方恒等）：}] 对任意 $\alpha\in\mathbb{R}$，有 $\sin^2\alpha+\cos^2\alpha=1$，它是三角学中最基本的恒等式。
		\item[\textbf{要点二（商数条件）：}] $\tan\alpha=\dfrac{\sin\alpha}{\cos\alpha}$ 仅在 $\cos\alpha\neq0$ 时成立，此时正切与正余弦可相互表达。
		\item[\textbf{要点三（变形方向）：}] 平方关系可直接改写为 $\sin^2\alpha=1-\cos^2\alpha$ 或 $\cos^2\alpha=1-\sin^2\alpha$，用于消元或降次。
	\end{description}

	\medskip
	\textbf{\textcolor{CExplanation}{几何本质（单位圆）}}

	在单位圆中，角 $\alpha$ 的终边与圆交于点 $P(\cos\alpha,\sin\alpha)$。由圆的方程 $x^2+y^2=1$ 立即得到 $\sin^2\alpha+\cos^2\alpha=1$。

	\medskip
	\textbf{\textcolor{CExplanation}{代数意义（参数约束）}}

	$\sin\alpha$ 与 $\cos\alpha$ 可视为满足 $u^2+v^2=1$ 的一对参数，因此 $|\sin\alpha|\le1$，$|\cos\alpha|\le1$，且两者取值相互制约。

	\medskip
	\textbf{\textcolor{CExplanation}{考点价值（考试指向）}}
	\begin{itemize}[leftmargin=*, itemsep=0.5em, topsep=0.4em]
		\item \textbf{考法一（知一求二）：} 已知 $\sin\alpha$ 或 $\cos\alpha$，通过平方关系求另一个，需注意象限符号。
		\item \textbf{考法二（恒等变形）：} 在化简或证明中，用 $\sin^2\alpha+\cos^2\alpha=1$ 代换 $1$，或用 $\sin^2\alpha=1-\cos^2\alpha$ 消去平方。
		\item \textbf{考法三（符号判定）：} 结合商数关系与象限，由 $\tan\alpha$ 符号判断 $\sin\alpha,\cos\alpha$ 同号或异号。
	\end{itemize}

	\medskip
	\textbf{\textcolor{CExplanation}{顿悟点}}

	把 $\sin^2\alpha+\cos^2\alpha=1$ 看作单位圆上的勾股定理，$\sin\alpha$ 和 $\cos\alpha$ 就是直角三角形的两条直角边（斜边为1）。

	\medskip
	\textbf{\textcolor{CExplanation}{使用场景}}
	\begin{itemize}[leftmargin=*, itemsep=0.5em, topsep=0.4em]
		\item \textbf{场景一（求值计算）：} 已知一个三角函数值及象限，求其它三角函数值。
		\item \textbf{场景二（化简表达式）：} 将 $\sin^2\theta$ 或 $\cos^2\theta$ 相互转化，简化复杂三角式。
		\item \textbf{场景三（解三角方程）：} 利用平方关系消元或引入辅助角。
	\end{itemize}
\end{explanationbox}
